\documentclass[11pt]{article}
\usepackage[margin=0.9in]{geometry}
\usepackage{amsmath,amssymb}
\usepackage{booktabs}
\usepackage{array}
\usepackage{tabularx}
\usepackage{xcolor}
\usepackage{multirow}
\usepackage[colorlinks=true,linkcolor=blue,urlcolor=blue]{hyperref}
\usepackage{enumitem}

\newcommand{\B}[1]{\textbf{#1}}
% colour code: local (our runs) = green, published OGBench = orange
\definecolor{locgreen}{rgb}{0.00,0.50,0.20}
\definecolor{puborange}{rgb}{0.85,0.45,0.00}
\newcommand{\loc}[1]{\textcolor{locgreen}{#1}}
\newcommand{\pub}[1]{\textcolor{puborange}{#1}}
\newcommand{\win}[1]{\B{\loc{#1}}}   % local number that beats published SOTA
\newcommand{\nr}{\textcolor{gray}{\emph{not run}}}
\newcommand{\inflight}[1]{\loc{#1}\,\textcolor{gray}{$^{\circ}$}} % @400k, run still finishing

\title{\bf Is bilinear value factorization a new SOTA on OGBench?}

\begin{document}
\maketitle

\begin{abstract}
We swap the monolithic MLP value head of three OGBench reference agents
(\B{GCIQL}, \B{HIQL}, \B{GCIVL}) for a rank-$d$ bilinear factorization
$Q(s,a,g)=\phi(s,a)^{\!\top}\psi(g)/\sqrt{d}$ and evaluate on \B{cube-single},
\B{antmaze-medium}, and \B{scene} (state obs). On cube, bilinear beats the
published SOTA for \emph{all three} methods once wired correctly (GCIQL
\win{90.4} vs \pub{68}, GCIVL \win{79.2} vs \pub{53}, HIQL \win{41.3} vs
\pub{15}); on antmaze the win is method-relative only, and on scene bilinear
loses with every actor. The main failure mode is identified: the bilinear TD
loss must not become the trainer of a representation the policy shares;
HIQL recovers exactly when $\psi$ is routed around its subgoal representation
$\zeta$.
\end{abstract}

\section{The factorization, per algorithm}
\label{sec:setup}
Offline GCRL setting per OGBench~\cite{park2025ogbench}: reward-free dataset
$\mathcal{D}$, hindsight-relabelled goals $g\sim p(g\mid\tau)$ (current / future
/ random mixture, OGBench defaults), sparse reward
$r(s,g)=\mathbf{1}[s\in S_g]-1$, termination mask $m=\mathbf{1}[s\notin S_g]$.
All value losses use the expectile loss~\cite{kostrikov2022iql}
$\ell^2_\tau(x)=|\tau-\mathbf{1}[x<0]|\,x^2$, $\tau\in[0.5,1)$; targets $\bar
V,\bar Q$; 2-ensembles throughout.\footnote{OGBench stabilizes the expectile
regression with a double-DQN-style split (expectile weight from the target-net
advantage, regressed difference from the online net); we keep it verbatim in
every variant.} We replace each value MLP with a rank-$d$ ($d{=}50$) two-tower
inner product, leaving losses, targets, and subgoal machinery untouched.

\B{GCIVL} (action-free goal-conditioned IQL~\cite{park2023hiql}; AWR~\cite{peng2019awr} actor):
\begin{equation}
\mathcal{L}_{V}^{\text{IVL}}
=\mathbb{E}\Bigl[\ell^2_\tau\bigl(r+\gamma m\,
\textstyle\min_i \bar V_i(s',g)-V(s,g)\bigr)\Bigr],
\qquad\text{factorized: }\;
V(s,g)=\tfrac{1}{\sqrt d}\,\phi(s)^{\!\top}\psi(g).
\label{eq:ivl}
\end{equation}

\B{GCIQL} (goal-conditioned IQL~\cite{kostrikov2022iql}; DDPG+BC~\cite{fujimoto2021td3bc} actor):
\begin{equation}
\mathcal{L}_{V}^{\text{IQL}}
=\mathbb{E}\Bigl[\ell^2_\tau\bigl(\textstyle\min_i\bar Q_i(s,a,g)-V(s,g)\bigr)\Bigr],
\qquad
\mathcal{L}_{Q}
=\mathbb{E}\Bigl[\textstyle\sum_{i}\bigl(Q_i(s,a,g)-r-\gamma m V(s',g)\bigr)^{2}\Bigr],
\label{eq:iql}
\end{equation}
factorized in \emph{both} heads:
$Q_i(s,a,g)=\tfrac{1}{\sqrt d}\phi_i(s,a)^{\!\top}\psi_i(g)$,
$V(s,g)=\tfrac{1}{\sqrt d}\phi_V(s)^{\!\top}\psi_V(g)$.

\B{HIQL}~\cite{park2023hiql} trains one value with Eq.~\eqref{eq:ivl} but reads
the goal through a learned, length-normalized subgoal representation
$\zeta:\mathcal{S}\times\mathcal{G}\to\mathbb{R}^k$ ($k{=}10$,
$\|\zeta\|=\sqrt k$):\footnote{The HIQL paper analyzes $V(s,\varphi(g))$ but
adopts the state-conditioned $\varphi([g,s])$ (our $\zeta(s,g)$) as the
empirically better variant; OGBench and we use exactly that.}
\begin{equation}
V(s,g)=f_\theta\bigl(s,\zeta(s,g)\bigr).
\label{eq:hiqlstock}
\end{equation}
$\zeta$ is HIQL's ``latent states as actions'': the high-level policy is trained
by AWR on $A^{h}=V(s_{t+n},g)-V(s_t,g)$ to output $z=\zeta(s_t,s_{t+n})$, and the
low-level policy consumes $z$. On state observations the low-level gradient into
$\zeta$ is off by default, so \B{the value loss is $\zeta$'s sole trainer}. This
makes the factorization ambiguous ($\psi$ may eat $g$ or $\zeta(s,g)$); we test
three forms:
\begin{align}
\text{(a) via shared rep (BVN-style~\cite{hong2023bvn}):}\quad
& V(s,g) = \tfrac{1}{\sqrt d}\,\phi(s)^{\!\top}\psi\bigl(\zeta(s,g)\bigr),
\label{eq:hiqlrep}\\
\text{(b) Schaul form~\cite{schaul2015uvfa}:}\quad
& V(s,g) = \tfrac{1}{\sqrt d}\,\phi(s)^{\!\top}\psi(g),
\label{eq:hiqlschaul}\\
\text{(c) two-head split:}\quad
& V = \text{stock Eq.~\eqref{eq:hiqlstock}}, \quad
V^{h}(s,g) = \tfrac{1}{\sqrt d}\,\phi(s)^{\!\top}
\psi(g)\big|_{\,\mathrm{sg}(\zeta)}.
\label{eq:hiqlsplit}
\end{align}
In (a) the bilinear loss \emph{replaces} the stock loss as $\zeta$'s trainer; in
(b) $\zeta$ gets no value gradient (the low-level gradient is re-enabled to keep
it trained); in (c) the stock $V$ keeps TD, low level, and $\zeta$-training,
while a stop-gradient bilinear head serves only the high-level advantage. The
GCIQL/GCIVL factorizations are Schaul-form by construction; HIQL's (a) is the
exception, and it matters (\S\ref{sec:hiqlmech}).

Besides each method's stock actor we test a distilled \emph{flow} actor
(one-step flow-matching BC + $Q$-max distillation, \`a la
FQL~\cite{park2025fql}); the actor turns out to be a first-order confound, so it
gets separate rows.

\section{What we evaluated}
\label{sec:results}
Table~\ref{tab:main} is the evaluation record; protocol in the caption. Cube:
bilinear beats the published SOTA for all three methods: GCIQL flow
\win{90.4} (\loc{93.1} @1M; $d{=}128/512$: 94.2/94.6; a \emph{domain} SOTA,
best published across all eight OGBench methods is \pub{68}), GCIVL \win{67.5}
(AWR) / \win{79.2} (flow), HIQL \win{32.8}--\win{41.3} via forms
(b)/(c). Rank is not the knob: $d\in\{50,128,512\}$ is flat and $d{=}5$
cliffs to 12--15. Antmaze: \win{83.7} vs \pub{71} is \emph{method-relative}
(published HIQL owns the domain at \pub{96}). Scene: bilinear loses with
\emph{both} actors: the one domain where the factorized value itself fails.
Which actor unlocks bilinear flips by domain, so any claim must name
(method, actor, domain):
\begin{center}\small
\begin{tabular}{@{}l ccc@{}}
\toprule
GCIQL bilinear $d{=}50$ & flow actor & DDPG+BC actor & stock MLP (best) \\
\midrule
cube-single  & \win{$90.4$} & \loc{$51.3$}\,$^{\alpha}$ & \loc{$71.7$} \\
antmaze-med  & \loc{$23.9$} & \win{$83.7$} & \loc{$69.5$} \\
scene        & \loc{$29.2$} & \loc{$25.4$} & \loc{$51.6$} \\
\bottomrule
\end{tabular}
\end{center}

\begin{table}[h]
\centering
\small
\setlength{\tabcolsep}{4.5pt}
\begin{tabular}{@{}l l l ccc@{}}
\toprule
Method & Value head & Actor & cube-single & antmaze-med & scene \\
\midrule
\multirow{6}{*}{\B{GCIQL}}
 & \pub{OGBench (pub.)} & native      & \pub{$68\pm6$}   & \pub{$71\pm4$}   & \pub{$51\pm4$} \\
 & MLP (native)         & DDPG+BC     & \loc{$70.3\pm4.7$}& \loc{$69.5\pm4.4$}& \loc{$51.6\pm3.6$} \\
 & MLP (native)         & flow        & \loc{$71.7\pm4.9$}& \loc{$33.3\pm5.2$}& \nr \\
 & \B{bilinear} $d{=}50$& flow        & \win{$90.4\pm3.4$}& \loc{$23.9\pm1.6$}& \loc{$29.2\pm3.2$} \\
 & \B{bilinear} $d{=}50$& DDPG+BC     & \loc{$51.3\pm8.0$}\,$^{\alpha}$& \win{$83.7\pm0.7$}& \loc{$25.4\pm0.3$}\,$^{s}$ \\
\midrule
\multirow{6}{*}{\B{HIQL}}
 & \pub{OGBench (pub.)} & native          & \pub{$15\pm3$}   & \pub{$96\pm1$}   & \pub{$38\pm3$} \\
 & MLP (native)         & hier.\ AWR      & \loc{$16.7\pm5.8$}& \loc{$95.3\pm0.6$}\,$^{m}$ & \nr \\
 & \B{bilinear} Eq.~\eqref{eq:hiqlrep} & hier.\ AWR & \loc{$18.5\pm0.2$}& \loc{$95.9\pm0.5$}\,$^{m}$ & \nr \\
 & \B{bilinear} Eq.~\eqref{eq:hiqlschaul} & hier.\ AWR & \win{$32.8\pm0.4$}& \nr & \nr \\
 & \B{bilinear} Eq.~\eqref{eq:hiqlsplit} & hier.\ AWR & \win{\inflight{$34.3\pm5.4$}}& \nr & \nr \\
 & \B{bilinear} Eq.~\eqref{eq:hiqlsplit} & flow       & \win{\inflight{$41.3\pm2.3$}}& \nr & \nr \\
\midrule
\multirow{5}{*}{\B{GCIVL}}
 & \pub{OGBench (pub.)} & native      & \pub{$53\pm4$}   & \pub{$72\pm8$}   & \pub{$42\pm4$} \\
 & MLP (native)         & AWR         & \loc{$49.6\pm2.0$}& \nr & \nr \\
 & MLP (native)         & flow        & \loc{$59.1\pm4.2$}& \nr & \nr \\
 & \B{bilinear} $d{=}50$& AWR         & \win{$67.5\pm8.7$}& \nr & \nr \\
 & \B{bilinear} $d{=}50$& flow        & \win{$79.2\pm0.4$}\,$^\dagger$& \nr & \nr \\
\bottomrule
\end{tabular}
\caption{\B{Does bilinear factorization beat SOTA?} State obs;
\texttt{evaluation/overall\_success} (\%) over 5 eval goals, mean\,$\pm$\,SD,
$n{=}3$ seeds @500k unless noted. \loc{Green}=local, \pub{orange}=OGBench
paper (8 seeds @1M); \B{bold green} = beats the published number for that
method/domain. $^{m}$matched protocol: 8 seeds, 1M, avg last-3 evals.
$^{\circ}$read @400k, runs finishing. $^{\alpha}$$\alpha{=}0.3$ vs stock
$\alpha{=}1.0$ (confounded, rerun owed). $^{s}$$n{=}2$ (third seed @300k:
20.8). $^\dagger$healthy batch; an earlier batch died at init and had
contaminated this cell as ``$47.5\pm38.8$''. HIQL cube cells use the tuned
\texttt{subgoal\_steps}${=}10$. \nr\ = not run.}
\label{tab:main}
\end{table}

\section{Why: the shared-representation mechanism}
\label{sec:hiqlmech}
The same factorization that buys GCIQL $+22$ and GCIVL $+18$--$30$ on cube moved
HIQL by $+1.8$ and \emph{hurt} its flow variant. The reason is in the algebra of
\S\ref{sec:setup}: $\zeta$ is at once the value's goal encoder, the high-level
action space, and the low-level goal input, and form~\eqref{eq:hiqlrep}
silently swaps its only training signal from the stock TD loss to the bilinear
one. GCIQL/GCIVL have no such shared module. A linear probe of $\zeta$ @500k
confirms: no rank collapse (participation ratio rises), but global geometry
degrades under~\eqref{eq:hiqlrep} (goal-xyz $R^2$ on random pairs
$0.14\!\to\!0.08$, rep-vs-goal-distance Spearman $0.115\!\to\!0.06$) while
subgoal-scale structure is spared. Cutting the value$\to\zeta$ path,
by~\eqref{eq:hiqlschaul} or~\eqref{eq:hiqlsplit}, is the fix
(Table~\ref{tab:hiqlmech}): every cell where the bilinear loss trains $\zeta$ is
flat-to-worse and capacity ($k\,10\!\to\!50$) does not rescue it; every cell
where it does not is a $2$--$2.7\times$ win over the published \pub{15}, and the
flow cell flips from bilinear's worst (14.5) to its best (41.3$^{\circ}$). The
matched-protocol antmaze control (MLP \loc{95.3} vs form-(a) \loc{95.9}, 8
seeds @1M) shows the warp costs nothing at that domain's ceiling, which is why
it was invisible where HIQL shines. The design rule: \B{the bilinear head must
not train a representation the policy shares}; when in doubt, give it a
stop-gradient path. Scene, where bilinear fails for GCIQL with both actors,
plausibly needs the value to compose multiple object-to-object interactions,
the one setting where a rank-50 product of a state tower and a goal tower is the
wrong bias; the~\eqref{eq:hiqlschaul}-vs-\eqref{eq:hiqlrep} axis is untested
there and is the obvious next probe.

\begin{table}[h]
\centering\small
\begin{tabular}{@{}l l l c@{}}
\toprule
Cell & who trains $\zeta$ & high-level advantage from & success \\
\midrule
MLP (baseline), Eq.~\eqref{eq:hiqlstock}\footnotemark & MLP value & MLP $V$ & \loc{$16.7\pm5.8$} \\
bilinear via shared rep, Eq.~\eqref{eq:hiqlrep} & \B{bilinear value} & bilinear $V$ & \loc{$18.5\pm0.2$} \\
\quad + $\zeta$ dim $k\!:10\!\to\!50$ & \B{bilinear value} & bilinear $V$   & \inflight{$24.3\pm1.0$} \\
bilinear, Schaul $\psi(g)$, Eq.~\eqref{eq:hiqlschaul} & low actor & bilinear $V$ & \win{$32.8\pm0.4$} \\
two-head split, Eq.~\eqref{eq:hiqlsplit} & MLP value & stop-grad bilinear head & \win{\inflight{$34.3\pm5.4$}} \\
\midrule
flow actors: MLP          & MLP value             & MLP $V$        & \loc{$22.7\pm5.2$} \\
flow, bilinear via shared rep & \B{bilinear value} & bilinear $V$  & \loc{$14.5\pm2.6$} \\
flow, two-head split      & MLP value             & stop-grad bilinear head & \win{\inflight{$41.3\pm2.3$}} \\
\bottomrule
\end{tabular}
\caption{\B{HIQL mechanism ladder} (cube-single, \texttt{subgoal\_steps}=10,
$d{=}50$, 3 seeds @500k; $^{\circ}$=@400k, finishing). \B{Bold ``who trains
$\zeta$''} = the bilinear TD loss shapes the shared subgoal representation.}
\label{tab:hiqlmech}
\end{table}
\footnotetext{Earlier drafts ran \texttt{subgoal\_steps}${=}25$ (MLP 24.6,
bilinear 15.9) and concluded ``bilinear doesn't help HIQL''; at the tuned
sg${=}10$ the baseline reproduces the published \pub{15} and that conclusion is
retired.}

\paragraph{Caveats.}
\begin{itemize}[nosep,leftmargin=1.4em]
\item Local = 500k/3 seeds vs published 1M/8 seeds; where both exist 500k
\emph{under}states bilinear ($90.4\!\to\!93.1$), but the winning cells still owe
a matched 8-seed 1M confirmation.
\item $^{\circ}$ cells are @400k with runs finishing; the ordering is already
fixed by the fully-final Schaul cell (32.8 @500k).
\item Confounds: bilinear $d{=}50$ has $2.06\times$ the value-head params of the
MLP (control owed); cube DDPG+BC ran $\alpha{=}0.3$ ($^{\alpha}$); scene is
untested on the~\eqref{eq:hiqlschaul}-vs-\eqref{eq:hiqlrep} axis.
\item Not run: GCIVL beyond cube, HIQL mechanism cells beyond cube, pixels.
\end{itemize}

\B{Provenance.} wandb \texttt{amsks/factored-fb} (pulled 2026-07-06, refreshed
from run-local \texttt{eval.csv} 2026-07-07), metric
\texttt{evaluation/overall\_success}; published numbers from
OGBench~\cite{park2025ogbench}, Table~2.

\begin{thebibliography}{9}
\bibitem{park2025ogbench}
S.~Park, K.~Frans, B.~Eysenbach, and S.~Levine.
\newblock OGBench: Benchmarking offline goal-conditioned RL.
\newblock In \emph{ICLR}, 2025. arXiv:2410.20092.

\bibitem{park2023hiql}
S.~Park, D.~Ghosh, B.~Eysenbach, and S.~Levine.
\newblock HIQL: Offline goal-conditioned RL with latent states as actions.
\newblock In \emph{NeurIPS}, 2023. arXiv:2307.11949.

\bibitem{kostrikov2022iql}
I.~Kostrikov, A.~Nair, and S.~Levine.
\newblock Offline reinforcement learning with implicit Q-learning.
\newblock In \emph{ICLR}, 2022. arXiv:2110.06169.

\bibitem{peng2019awr}
X.~B. Peng, A.~Kumar, G.~Zhang, and S.~Levine.
\newblock Advantage-weighted regression: Simple and scalable off-policy
reinforcement learning.
\newblock arXiv:1910.00177, 2019.

\bibitem{fujimoto2021td3bc}
S.~Fujimoto and S.~S. Gu.
\newblock A minimalist approach to offline reinforcement learning.
\newblock In \emph{NeurIPS}, 2021. arXiv:2106.06860.

\bibitem{schaul2015uvfa}
T.~Schaul, D.~Horgan, K.~Gregor, and D.~Silver.
\newblock Universal value function approximators.
\newblock In \emph{ICML}, 2015.

\bibitem{hong2023bvn}
Z.-W. Hong, G.~Yang, and P.~Agrawal.
\newblock Bilinear value networks.
\newblock In \emph{ICLR}, 2023. arXiv:2204.13695.

\bibitem{park2025fql}
S.~Park, Q.~Li, and S.~Levine.
\newblock Flow Q-learning.
\newblock In \emph{ICML}, 2025. arXiv:2502.02538.
\end{thebibliography}

\end{document}
